% !TEX root = ../proyect.tex

\chapter{Implementación del procedimiento de análisis}\label{cap:implementacion-procedimiento}

\section{Estructura general del procedimiento}\label{sec:estructura-procedimiento}

El procedimiento de análisis se implementa como una secuencia de fases reproducibles. Cada fase recibe una entrada concreta, ejecuta un conjunto delimitado de operaciones y exporta los resultados en formatos CSV o JSON, así como en elementos visuales como imágenes, gráficos o tablas. Esta organización permite repetir partes del estudio sin rehacer todo el proyecto y facilita que las cifras, tablas y figuras usadas posteriormente en la memoria puedan rastrearse hasta un script y un fichero de salida.

La cadena principal trabaja sobre datos de \texttt{BTCUSDT} con frecuencia de cinco minutos. Primero se obtiene y valida el dataset limpio; después se construyen retornos, medidas de volatilidad realizada y objetivos futuros; a continuación se aplican herramientas exploratorias, contrastes estadísticos, análisis lineal, reconstrucción del espacio de estados y comparación predictiva. Las fases finales incorporan una comprobación sintética con el mapa logístico y un modelo HAR-logRV compacto, exportable al MVP web.

La estructura del repositorio distingue entre datos brutos, datos procesados, resultados y código. Los ficheros originales se conservan en \texttt{data/\allowbreak raw/}; los datasets generados se guardan en \texttt{data/\allowbreak processed/}; las tablas, informes y figuras se escriben en \texttt{reports/}, \texttt{reports/\allowbreak tables/} y \texttt{reports/\allowbreak figures/}. Esta separación evita mezclar entradas, resultados intermedios y productos finales del análisis.

En este capítulo se describe la implementación del procedimiento, no la discusión completa de los resultados. Por ello, las secciones siguientes explican qué calcula cada fase, qué criterios se han fijado y qué artefactos produce. La interpretación conjunta de las salidas se reserva para el capítulo de resultados y discusión.

\section{Entorno de desarrollo y herramientas utilizadas}\label{sec:entorno-desarrollo}

El procedimiento técnico se implementa en Python. La decisión de usar scripts independientes por fase responde a una necesidad práctica: el análisis combina procesos de naturaleza distinta, desde validación de datos hasta cálculo de correlogramas, gráficos de recurrencia, reconstrucción y predicción. Separar cada bloque en un script reduce dependencias internas y permite revisar una fase sin alterar el resto de la cadena.

La construcción inicial del dataset limpio utiliza \texttt{pandas}, porque los datos proceden de ficheros mensuales comprimidos y es necesario leer, concatenar, ordenar y exportar grandes tablas de velas. En cambio, varias fases posteriores se apoyan en módulos propios y, cuando es viable, en biblioteca estándar. Por ejemplo, la validación integral del CSV se implementa sin \texttt{pandas}; los gráficos SVG se generan mediante funciones propias; y los cálculos de ACF, PACF, periodograma, ADF y KPSS se implementan en módulos auxiliares del repositorio.

Los módulos de apoyo separan operaciones comunes. \texttt{data\_\allowbreak loading.py} centraliza lectura y escritura de CSV; \texttt{preprocessing.py} construye las variables de Fase 1; \texttt{volatility.py} contiene las sumas móviles pasadas y futuras; \texttt{stationarity.py} implementa herramientas de estacionariedad; \texttt{spectral.py} calcula correlogramas y espectro; \texttt{recurrence.py} genera gráficos de recurrencia; y \texttt{plotting.py} agrupa utilidades ligeras para figuras SVG.

La salida de cada fase se guarda en formatos simples. Las tablas se escriben en CSV, los resúmenes de parámetros en JSON cuando procede, las matrices o referencias numéricas en NPZ y las figuras en SVG o PNG. Esta elección simplifica la reutilización de resultados por la memoria y por el MVP web, sin depender de objetos internos de Python.

\section{Obtención y validación inicial de los datos}\label{sec:obtencion-datos}

El dato bruto procede de los ficheros mensuales de \textit{klines} de Binance Spot para el par \texttt{BTCUSDT} con intervalo de cinco minutos. Los ficheros se almacenan en \texttt{data/\allowbreak raw/} con nombres del tipo \texttt{BTCUSDT-5m-2024-01.zip}. Cada ZIP contiene un CSV sin cabecera con las columnas originales de Binance.

El script \texttt{build\_\allowbreak btc\_\allowbreak 5m\_\allowbreak clean.py} lee todos los ZIP disponibles en \texttt{data/\allowbreak raw/}, extrae el CSV interno de cada archivo y concatena las tablas resultantes. Durante la limpieza se convierte \texttt{open\_\allowbreak time} a fecha UTC, se detecta automáticamente si la unidad temporal viene en milisegundos, microsegundos o nanosegundos, se convierten las columnas numéricas, se ordenan las filas por tiempo y se eliminan duplicados por \texttt{open\_\allowbreak time}. La salida se reduce a las columnas necesarias para el análisis: \texttt{open\_\allowbreak time}, \texttt{open}, \texttt{high}, \texttt{low}, \texttt{close}, \texttt{volume}, \texttt{trades} y \texttt{taker\_\allowbreak buy\_\allowbreak quote}.

La validación formal se realiza con \texttt{validate.py}. Este script comprueba que el CSV limpio tiene las columnas esperadas, que el rango temporal coincide con el intervalo declarado, que la frecuencia modal es de cinco minutos, que no existen diferencias temporales distintas de cinco minutos, que \texttt{open\_\allowbreak time} es estrictamente creciente, que no hay duplicados, huecos temporales, nulos ni errores de conversión numérica, y que los precios son positivos y coherentes con la relación OHLC.

La Fase 0 escribe varias tablas auxiliares en \texttt{reports/\allowbreak tables/}: resumen de calidad, frecuencia temporal, nulos, comprobaciones numéricas y reporte de huecos. En la ejecución documentada, el dataset limpio contiene 245.088 observaciones, desde el 1 de enero de 2024 a las 00:00 hasta el 30 de abril de 2026 a las 23:55. Esta fase fija la base empírica sobre la que se construyen las variables del resto del trabajo.

\section{Construcción de variables y medidas de volatilidad}\label{sec:construccion-variables}

La Fase 1 transforma el CSV limpio en el dataset de trabajo \texttt{data/\allowbreak processed/\allowbreak btc\_\allowbreak 5m\_\allowbreak features.csv}. El script principal es \texttt{build\_\allowbreak phase1\_\allowbreak features.py}, que delega la construcción de variables en \texttt{preprocessing.py} y las sumas móviles de volatilidad en \texttt{volatility.py}.

Primero se calcula el logaritmo del cierre, \texttt{log\_\allowbreak close}, y a partir de él el retorno logarítmico de cinco minutos,
\[
    r_t = \log(C_t) - \log(C_{t-1}).
\]
Sobre el retorno se obtienen \texttt{r2}, \texttt{abs\_\allowbreak r}, el rango logarítmico \texttt{hl\_\allowbreak range}, y las transformaciones \texttt{log\_\allowbreak volume} y \texttt{log\_\allowbreak trades}. Estas variables permiten conservar información de precio, intensidad de movimiento y actividad de mercado.

La volatilidad realizada se implementa como suma móvil de retornos al cuadrado. Las variables \texttt{rv\_\allowbreak past\_\allowbreak 12}, \texttt{rv\_\allowbreak past\_\allowbreak 48} y \texttt{rv\_\allowbreak past\_\allowbreak 288} agregan, respectivamente, 1 hora, 4 horas y 1 día de datos pasados. La convención es inclusiva hasta \(t\): la ventana pasada de tamaño \(h\) usa los retornos desde \(t-h+1\) hasta \(t\).

Los objetivos futuros se construyen con una convención estrictamente posterior. \texttt{rv\_\allowbreak future\_\allowbreak 12} y \texttt{rv\_\allowbreak future\_\allowbreak 48} agregan retornos desde \(t+1\) hasta el horizonte correspondiente. Esta definición evita que el retorno del instante \(t\) entre simultáneamente en las variables explicativas y en el objetivo futuro.

Sobre todas las volatilidades realizadas se aplica una transformación logarítmica con \(\varepsilon=10^{-12}\). Así se obtienen, entre otras, \texttt{log\_\allowbreak rv\_\allowbreak past\_\allowbreak 12} y \texttt{log\_\allowbreak rv\_\allowbreak future\_\allowbreak 12}. Tras eliminar las filas iniciales sin histórico suficiente y las filas finales sin futuro disponible, el dataset procesado queda con 244.752 observaciones, desde el 2 de enero de 2024 a las 00:00 hasta el 30 de abril de 2026 a las 19:55.

La serie central del procedimiento queda definida como
\[
    v_t = \texttt{log\_\allowbreak rv\_\allowbreak past\_\allowbreak 12},
\]
mientras que el objetivo predictivo principal es \texttt{log\_\allowbreak rv\_\allowbreak future\_\allowbreak 12}. Esta elección no se discute aquí como resultado final; se fija como salida técnica de la construcción de variables y se evalúa en las fases posteriores.

\section{Visualización inicial y estadísticos descriptivos}\label{sec:visualizacion-descriptivos}

La Fase 2 aplica una primera batería de herramientas descriptivas sobre \texttt{btc\_\allowbreak 5m\_\allowbreak features.csv}. El script \texttt{phase2\_\allowbreak general\_\allowbreak tools.py} lee el dataset procesado, convierte las columnas numéricas necesarias y genera figuras temporales, estadísticos descriptivos y una tabla de eventos extremos.

Las figuras generadas cubren el precio de cierre, el logaritmo del precio, los retornos logarítmicos, los retornos absolutos, \texttt{log\_\allowbreak rv\_\allowbreak past\_\allowbreak 12}, \texttt{log\_\allowbreak rv\_\allowbreak future\_\allowbreak 12}, volumen y número de operaciones. Las salidas se escriben como SVG en \texttt{reports/\allowbreak figures/}. El objetivo de estas figuras es inspeccionar la escala de las variables, localizar episodios extremos y comprobar visualmente la diferencia entre precio, retorno e intensidad de movimiento.

La fase también calcula estadísticos descriptivos para las variables principales: media, desviación típica, mínimo, percentiles, máximo, asimetría y exceso de curtosis. Además, identifica eventos extremos como el retorno más negativo, el retorno más positivo, el mayor retorno absoluto, la mayor volatilidad realizada pasada, la mayor volatilidad realizada futura, el mayor volumen y el mayor número de operaciones.

Esta fase no realiza contrastes formales ni extrae conclusiones sobre caos o predictibilidad. Su función dentro del procedimiento es preparar una lectura inicial del dataset procesado y documentar qué variables merecen análisis más detallado en las fases de estacionariedad, correlación, recurrencia y reconstrucción.

\section{Análisis de estacionariedad}\label{sec:analisis-estacionariedad}

La Fase 3 estudia estacionariedad y transformaciones sobre cinco series: \texttt{close}, \texttt{log\_\allowbreak close}, \texttt{r}, \texttt{abs\_\allowbreak r} y \texttt{log\_\allowbreak rv\_\allowbreak past\_\allowbreak 12}. El script \texttt{phase3\_\allowbreak stationarity.py} genera gráficos temporales, gráficos de media y desviación típica móvil, contrastes ADF y contrastes KPSS.

La ventana móvil se fija en 288 observaciones, equivalente a un día de datos de cinco minutos. Para el ADF se prueban retardos candidatos \(\{0,1,2,3,6,12\}\) y se selecciona el retardo mediante BIC. El KPSS se implementa con constante y varianza de largo plazo tipo Newey-West. En esta fase no se usa \texttt{statsmodels}; los módulos propios devuelven estadísticos y rangos aproximados de \textit{p-value} usando valores críticos asintóticos habituales.

Las salidas de la fase incluyen figuras temporales por serie, figuras de momentos móviles y tablas CSV con resultados ADF, selección de retardos ADF, resultados KPSS y resumen de momentos móviles. Estas salidas permiten justificar qué transformaciones son adecuadas para continuar el análisis, pero la lectura completa de estacionariedad y cambios de régimen se desarrolla en el capítulo de resultados.

\section{Análisis de autocorrelación y espectro}\label{sec:autocorrelacion-espectro}

La Fase 4 analiza dependencia lineal y estructura espectral en tres series: \texttt{r}, \texttt{abs\_\allowbreak r} y \texttt{log\_\allowbreak rv\_\allowbreak past\_\allowbreak 12}. El script \texttt{phase4\_\allowbreak correlogram\_\allowbreak spectrum.py} calcula ACF, PACF y periodograma para cada una de ellas.

La ACF se calcula hasta 2016 retardos, equivalente a una semana de datos de cinco minutos. También se genera una vista específica hasta 288 retardos, equivalente a un día. La PACF se calcula hasta 288 retardos mediante recursiones de Levinson-Durbin. El periodograma se obtiene tras restar la media y se representa frente al periodo equivalente en retardos de cinco minutos, utilizando escala logarítmica para poder comparar escalas de una hora, un día y una semana.

La fase guarda tablas con valores de ACF, valores de PACF, picos espectrales, potencia en periodos de referencia y resumen de correlograma. Las figuras correspondientes se escriben en SVG. Las bandas de ACF/PACF se tratan como referencias visuales, no como contraste exacto de ruido blanco, porque el tamaño muestral, las colas pesadas y la heterocedasticidad pueden hacer que retardos de magnitud pequeña resulten formalmente significativos.

El propósito de esta fase es documentar la diferencia entre retorno firmado e intensidad de movimiento. La interpretación final de los retardos relevantes y de los picos espectrales se reserva al capítulo de resultados; en la implementación basta con fijar el protocolo de cálculo y las salidas generadas.

\section{Gráficos de recurrencia}\label{sec:graficos-recurrencia}

La Fase 5 introduce gráficos de recurrencia iniciales sobre ventanas representativas. No se calcula una matriz de recurrencia completa sobre toda la muestra, porque el tamaño del dataset haría inviable una matriz \(N \times N\) y, además, su interpretación visual sería poco útil. En su lugar se seleccionan ventanas locales de 2.000 observaciones.

El script \texttt{phase5\_\allowbreak initial\_\allowbreak recurrence.py} selecciona cuatro ventanas: una ventana tranquila, una ventana de alta volatilidad, una ventana reciente y una ventana central representativa. La ventana tranquila se elige mediante medias móviles de \texttt{log\_\allowbreak rv\_\allowbreak past\_\allowbreak 12}; la de alta volatilidad se centra en el máximo de esa misma serie; la reciente toma las últimas observaciones disponibles; y la central se sitúa en la mitad del dataset procesado.

En cada ventana se analizan tres series: \texttt{r}, \texttt{abs\_\allowbreak r} y \texttt{log\_\allowbreak rv\_\allowbreak past\_\allowbreak 12}. Cada fragmento se normaliza mediante z-score calculado dentro de la propia ventana. La distancia utilizada es absoluta en una dimensión y el umbral \(\varepsilon\) se ajusta para alcanzar una tasa de recurrencia objetivo del 5\%. Los gráficos se guardan en PNG, donde el píxel negro representa un par recurrente y el píxel blanco un par no recurrente.

La fase escribe dos tablas principales: \texttt{phase5\_\allowbreak selected\_\allowbreak windows.csv}, con la descripción y límites temporales de las ventanas, y \texttt{phase5\_\allowbreak recurrence\_\allowbreak parameters.csv}, con los parámetros usados para cada serie y ventana. En este punto los gráficos de recurrencia se usan como herramienta exploratoria sobre series escalares; no constituyen todavía una reconstrucción del espacio de estados ni una prueba de caos. Su papel es preparar visualmente la transición hacia las fases posteriores de filtrado, contraste y reconstrucción.

\section{Filtrado lineal mediante modelos autorregresivos}\label{sec:filtrado-lineal}

La Fase 6 introduce un filtrado lineal sobre la serie principal del estudio, \(v_t = \texttt{log\_\allowbreak rv\_\allowbreak past\_\allowbreak 12}\). Antes del ajuste, la serie se estandariza utilizando exclusivamente la media y la desviación típica del tramo de entrenamiento. Esta decisión evita que información posterior al corte temporal intervenga en la escala usada por el modelo.

El script \texttt{phase6\_\allowbreak linear\_\allowbreak filtering.py} fija el final del entrenamiento en \texttt{2025-06-30 23:55:00}. Sobre ese tramo ajusta modelos AR(\(p\)) para \(p=1,\ldots,100\), mediante ecuaciones de Yule-Walker. Para cada orden se calcula la varianza de innovación, AIC y BIC. El orden seleccionado es el que minimiza BIC; en la ejecución documentada se obtiene un AR(49).

Una vez seleccionado el orden, el modelo se aplica sobre la serie completa estandarizada para obtener valores ajustados \textit{one-step} y residuos. Estos residuos se guardan en \texttt{reports/\allowbreak tables/\allowbreak phase6\_\allowbreak residual\_\allowbreak series.csv}, junto con el valor original de \texttt{log\_\allowbreak rv\_\allowbreak past\_\allowbreak 12}, su versión estandarizada, el ajuste AR y el residuo correspondiente. Esta tabla constituye la entrada principal de la fase siguiente.

La fase genera también tablas de selección del orden autorregresivo, coeficientes del AR seleccionado, estadísticos descriptivos de residuos, valores ACF/PACF de residuos, contrastes Ljung-Box y comparación de correlogramas entre la serie estandarizada y los residuos. Las figuras asociadas incluyen la serie temporal de residuos, su histograma, ACF/PACF y gráficos de recurrencia sobre las mismas ventanas seleccionadas en la Fase 5.

El papel de esta fase es doble. Por un lado, proporciona un benchmark lineal fuerte para las fases predictivas. Por otro, permite separar la dependencia lineal en media de la estructura que pueda permanecer en los residuos. La interpretación de si esa estructura residual es relevante se pospone a los contrastes de la Fase 7 y a la discusión de resultados.

\section{Contrastes de no linealidad}\label{sec:contrastes-no-linealidad}

La Fase 7 aplica contrastes de dependencia residual sobre dos series alineadas: la serie principal estandarizada, \texttt{z\_\allowbreak log\_\allowbreak rv\_\allowbreak past\_\allowbreak 12}, y los residuos del AR(49). El objetivo de la fase no es demostrar caos, sino comprobar si, tras retirar una componente lineal en media, siguen apareciendo dependencias compatibles con heterocedasticidad, estructura temporal en varianza o desviaciones frente a una secuencia i.i.d.

El script \texttt{phase7\_\allowbreak nonlinearity\_\allowbreak tests.py} toma como entrada \texttt{phase6\_\allowbreak residual\_\allowbreak series.csv}. En primer lugar calcula contrastes Ljung-Box sobre los cuadrados de ambas series, con retardos 12, 24, 48, 96, 288 y 2016. Esta prueba se usa para evaluar dependencia en segundo momento, especialmente relevante en series financieras donde la volatilidad tiende a agruparse.

A continuación se aplica un contraste ARCH-LM aproximado sobre los residuos del AR. La implementación evita construir una regresión masiva con matriz de diseño explícita y usa una aproximación mediante Yule-Walker sobre residuos cuadrados centrados. Los retardos considerados son 12, 24, 48, 96 y 288. La salida se guarda en \texttt{phase7\_\allowbreak arch\_\allowbreak lm.csv}, junto con el estadístico, el \(p\)-valor y una nota sobre la aproximación empleada.

La fase incorpora también el contraste BDS en ventanas representativas. Para mantener el coste computacional controlado, se trabaja sobre ventanas de 3.000 observaciones, reutilizando los centros de las ventanas seleccionadas previamente cuando están disponibles. Se evalúan dimensiones \(m=2,3,4\) y umbrales \(\epsilon\) iguales a 0,5, 1,0 y 1,5 desviaciones típicas, dado que cada ventana se normaliza antes del contraste. Los resultados se escriben en \texttt{phase7\_\allowbreak bds\_\allowbreak results.csv} y se resumen mediante una figura de \(p\)-valores.

Por último, se realiza una comparación frente a series barajadas. Para cada ventana y serie se calcula la autocorrelación de cuadrados en retardos 1, 12 y 288, y se compara con 50 permutaciones aleatorias generadas con semilla fija. Esta prueba conserva la distribución marginal de la ventana, pero destruye el orden temporal. La comparación permite distinguir entre efectos puramente distribucionales y dependencias asociadas al orden de la serie.

Las salidas de la Fase 7 incluyen tablas de ventanas, Ljung-Box sobre cuadrados, ARCH-LM, resultados BDS, resúmenes por ventana y comparación con barajado. Las figuras principales son la ACF de cuadrados de la serie estandarizada, la ACF de cuadrados de los residuos, el resumen BDS y la comparación de autocorrelaciones frente a permutaciones. La lectura metodológica es prudente: la fase busca evidencia de dependencia residual, no una prueba directa de dinámica caótica determinista.

\section{Reconstrucción del espacio de estados}\label{sec:implementacion-reconstruccion}

La Fase 8 implementa la reconstrucción del espacio de estados de la serie principal. La serie de partida vuelve a ser \(\texttt{log\_\allowbreak rv\_\allowbreak past\_\allowbreak 12}\), estandarizada con media y desviación típica calculadas solo en el tramo de entrenamiento. La reconstrucción se realiza mediante retardos temporales siguiendo la convención causal
\begin{equation}
    X_t = [z_t, z_{t-\tau}, z_{t-2\tau}, \ldots, z_{t-(m-1)\tau}].
\end{equation}
Esta orientación garantiza que cada vector reconstruido use únicamente información disponible en \(t\) o anterior a \(t\), requisito necesario para su uso posterior en predicción.

El script \texttt{phase8\_\allowbreak state\_\allowbreak space\_\allowbreak reconstruction.py} calcula primero el retardo \(\tau\) mediante información mutua media. Para ello discretiza la serie de entrenamiento en 32 contenedores definidos por cuantiles y evalúa la información mutua para retardos entre 1 y 288. La regla de selección busca el primer mínimo local de la curva AMI; si no aparece, contempla un criterio alternativo por codo. En la ejecución documentada el retardo seleccionado es \(\tau=137\). La fase calcula además una curva AMI para una versión barajada de la serie, usada como referencia visual.

La dimensión de embedding se estudia con falsos vecinos cercanos y con el método de Cao. El cálculo se realiza sobre una submuestra equiespaciada de 3.000 puntos del entrenamiento, con dimensión máxima 15. Para FNN se usan los criterios \(R_{\text{tol}}=10\) y \(A_{\text{tol}}=2\), y se fija como referencia un umbral del 5\%. El método de Cao calcula las secuencias \(E1\) y \(E2\) para evaluar estabilización de la dimensión. La regla final combina ambas fuentes: si FNN y Cao son compatibles, se adopta la mayor dimensión sugerida; si discrepan, se prioriza FNN y se conserva Cao como cautela metodológica. En la ejecución documentada se adopta \(m=5\).

Con los parámetros seleccionados se construyen embeddings separados para entrenamiento y prueba. La separación temporal es la misma que en las fases anteriores: el entrenamiento termina en \texttt{2025-06-30 23:55:00} y el tramo posterior queda reservado para evaluación. Los vectores se guardan en \texttt{data/\allowbreak processed/\allowbreak phase8\_\allowbreak embedding\_\allowbreak train.npz} y \texttt{data/\allowbreak processed/\allowbreak phase8\_\allowbreak embedding\_\allowbreak test.npz}, junto con índices y metadatos. Los parámetros finales se almacenan en \texttt{reports/\allowbreak tables/\allowbreak phase8\_\allowbreak selected\_\allowbreak embedding\_\allowbreak params.json}.

La fase genera tablas de AMI, AMI barajada, FNN, Cao y una muestra del embedding. También produce figuras de AMI, comparación original frente a barajada, FNN, Cao, proyecciones 2D/3D globales y proyecciones 2D por ventanas representativas. Estas visualizaciones no se interpretan como prueba de atractor; se usan para documentar la representación de estados que alimenta las fases de cuantificación y predicción local.

\section{Cuantificación de la dinámica reconstruida}\label{sec:cuantificacion-dinamica}

La Fase 9 cuantifica la dinámica asociada al embedding seleccionado. El script \texttt{phase9\_\allowbreak dynamics\_\allowbreak quantification.py} carga los parámetros de la Fase 8, reconstruye los vectores de entrenamiento y compara la serie original con una versión barajada reproducible. Esta comparación es importante porque permite distinguir entre propiedades que dependen del orden temporal y propiedades que podrían aparecer por la distribución marginal de la serie.

La primera familia de indicadores corresponde a la dimensión de correlación. La implementación toma una submuestra equiespaciada de 2.500 vectores y calcula distancias euclídeas entre pares válidos, excluyendo vecinos temporales mediante una ventana de Theiler igual a \(m\tau\). A partir de esas distancias se definen 32 radios entre cuantiles y se estima la curva \(C(r)\), su representación log-log y la pendiente local. La salida incluye tanto los valores de la curva como un resumen prudente de \(D_2\), que registra si existe o no una meseta estable.

La segunda familia de indicadores es la divergencia local de trayectorias cercanas mediante el método de Rosenstein. Para reducir el coste computacional, se toma un bloque continuo de 3.000 vectores situado en la zona central del entrenamiento. Se identifican vecinos no triviales, se calcula la evolución de la distancia media durante 60 pasos de cinco minutos y se ajusta una recta en el intervalo \(k=1,\ldots,30\). La pendiente se guarda tanto por paso de cinco minutos como reescalada por hora.

La tercera familia de indicadores es la entropía de permutación. Se calcula para órdenes 3, 4, 5, 6 y 7, usando dos retardos: 1 y \(\tau\). El mismo cálculo se repite sobre la serie barajada. Esta medida resume la diversidad de patrones ordinales y complementa los indicadores basados en distancias, ya que no requiere una interpretación geométrica directa del embedding.

La fase guarda tablas para la curva de dimensión de correlación, la curva de Rosenstein y la entropía de permutación. Además, escribe resúmenes JSON con los parámetros y resultados agregados. Las figuras asociadas muestran la curva log-log de correlación, la pendiente local, la divergencia de Rosenstein, la entropía de permutación y una comparación agregada entre original y barajada. La interpretación de estos indicadores se desarrolla en el capítulo de resultados; aquí quedan definidos los cálculos y las salidas generadas.

\section{Contrastes con barajado y datos subrogados}\label{sec:contrastes-barajado}

La Fase 10 amplía el control frente a series de referencia. La comparación con una serie barajada destruye todo orden temporal, pero no es suficiente para evaluar si ciertos resultados pueden explicarse por estructura lineal o por la distribución marginal. Por ello, el script \texttt{phase10\_\allowbreak surrogate\_\allowbreak tests.py} genera tres grupos de referencia: permutaciones aleatorias, series \textit{phase-randomized} y series AAFT.

La fase trabaja sobre una ventana de entrenamiento de 8.192 observaciones. Este tamaño es una potencia de dos, requisito práctico para la generación eficiente de subrogadas \textit{phase-randomized} mediante FFT. La ventana se toma del tramo de entrenamiento y se estandariza con los parámetros de train. Los parámetros \(\tau\), \(m\) y la convención de embedding proceden de la Fase 8, mientras que la exclusión temporal se fija de forma coherente con la ventana de Theiler usada en cuantificación.

Las series barajadas conservan los valores observados, pero eliminan su orden temporal. Las series \textit{phase-randomized} se generan aleatorizando fases en el dominio frecuencial, manteniendo de forma aproximada el espectro lineal. Las series AAFT aplican una transformación por rangos, aleatorización de fases y reordenación final para aproximar simultáneamente distribución marginal y estructura lineal. En la ejecución se generan 50 series barajadas, 39 \textit{phase-randomized} y 39 AAFT, todas con semilla fija para asegurar reproducibilidad.

Para cada réplica se calculan los mismos estadísticos principales: dimensión de correlación aproximada, pendiente de Rosenstein por paso y por hora, entropía de permutación con retardo 1 y entropía de permutación con retardo \(\tau\). Por coste computacional, esta fase usa una configuración más compacta que la Fase 9: muestra de 700 vectores para dimensión de correlación, 24 radios, bloque de 1.600 vectores para Rosenstein, 550 referencias, \(k_{\max}=40\) y ajuste en \(k=1,\ldots,20\).

Las salidas se guardan en \texttt{reports/\allowbreak tables/}. Incluyen el estadístico original, la configuración de la fase, las réplicas de cada grupo y un resumen por estadístico con media, desviación típica, percentiles, distancia estandarizada y \(p\)-valor empírico. Las figuras generadas muestran boxplots agregados, boxplots separados por tipo de estadístico, histogramas de Lyapunov y entropía, y ejemplos de serie original frente a subrogadas.

El objetivo de esta fase es someter los indicadores de la Fase 9 a referencias más exigentes. Si una diferencia frente a barajado desaparece frente a subrogadas que conservan espectro o distribución, la lectura debe limitarse. Por tanto, la Fase 10 actúa como mecanismo de control para que el capítulo de resultados pueda distinguir entre estructura temporal general, dependencia lineal, efectos distribucionales e indicios no lineales más robustos.

\section{Predicción local de volatilidad}\label{sec:prediccion-local}

La Fase 11 evalúa la capacidad predictiva del espacio reconstruido sobre el objetivo \texttt{log\_\allowbreak rv\_\allowbreak future\_\allowbreak 12}. La prueba se plantea como una evaluación empírica fuera de muestra. No se utiliza para demostrar caos, sino para comprobar si la representación de estados construida en la Fase 8 permite mejorar referencias predictivas simples bajo una separación temporal estricta.

El script \texttt{phase11\_\allowbreak local\_\allowbreak state\_\allowbreak space\_\allowbreak prediction.py} carga los parámetros seleccionados en la reconstrucción, en particular \(\tau=137\) y \(m=5\), y prepara tres particiones cronológicas: entrenamiento hasta el 30 de junio de 2025, validación hasta el 31 de diciembre de 2025 y test desde el 1 de enero de 2026. La serie de entrada es \texttt{log\_\allowbreak rv\_\allowbreak past\_\allowbreak 12}; el objetivo es \texttt{log\_\allowbreak rv\_\allowbreak future\_\allowbreak 12}; y el horizonte es de 12 velas, equivalente a una hora.

Los modelos comparados son media histórica, persistencia, AR(49) recursivo a horizonte 12, vecino más próximo y media de \(k\) vecinos. Para el método local se construyen vectores de embedding en cada partición y se buscan vecinos mediante un KD-tree exacto. La rejilla inicial de \(k\) es
\[
    k \in \{2,3,5,10,20,50\}.
\]
La selección de \(k\) se realiza exclusivamente en validación, usando RMSE como criterio. El conjunto de test no interviene en esta selección.

La fase incorpora varias reglas anti-\textit{leakage}. En validación, los vecinos se buscan solo dentro de entrenamiento. En test, la biblioteca de vecinos permitida es entrenamiento más validación. Además, un candidato solo es elegible si su etiqueta futura ya sería conocida en el instante de consulta, es decir, si cumple \(\texttt{candidate\_\allowbreak index}+12 \leq \texttt{query\_\allowbreak index}\). También se aplica una ventana de Theiler de \(\max(\texttt{theiler},m\tau)\), que en esta configuración queda en 685 retardos, para evitar vecinos temporalmente triviales.

Las salidas de la fase se escriben en \texttt{reports/\allowbreak tables/}. Incluyen el resumen de particiones, la selección de \(k\) en validación, las métricas de test, una muestra de predicciones fuera de muestra y un resumen JSON con los parámetros, controles de fuga de información y métricas agregadas. Las figuras generadas muestran la curva de validación por \(k\), la comparación real frente a predicho en una ventana de test, la evolución de errores, el histograma de errores y la comparación de métricas entre modelos.

\section{Experimentos de robustez y configuraciones alternativas}\label{sec:experimentos-robustez}

La Fase 12 extiende la evaluación predictiva de la Fase 11 sin introducir una nueva familia de modelos. Su objetivo es comprobar si el comportamiento del predictor local depende de una rejilla de \(k\) demasiado limitada o de la elección operativa \(m=5\), frente a la dimensión más alta sugerida por el método de Cao.

El script \texttt{phase12\_\allowbreak prediction\_\allowbreak sensitivity.py} evalúa dos configuraciones:
\[
    (\tau,m)=(137,5)
    \quad\text{y}\quad
    (\tau,m)=(137,14).
\]
En la primera, la memoria efectiva del embedding es \((5-1)137=548\) velas, aproximadamente 45,7 horas. En la segunda, la memoria efectiva es \((14-1)137=1781\) velas, aproximadamente 148,4 horas. La ventana de Theiler se fija como \(\tau m\) en cada configuración.

La rejilla de vecinos se amplía a
\[
    k \in \{2,3,5,10,20,50,100,200\}.
\]
La configuración \(m=5\) se evalúa sobre 5.000 puntos equiespaciados de validación y test. La configuración \(m=14\) se evalúa sobre 1.000 puntos, debido al mayor coste computacional asociado a la dimensión alta y a la menor densidad local del espacio reconstruido.

La lógica de evaluación mantiene las mismas restricciones temporales que la Fase 11: validación busca vecinos solo en entrenamiento; test busca vecinos en entrenamiento más validación; los candidatos deben tener etiqueta futura disponible; y test no participa en la selección de \(k\). Además de kNN, se conservan las referencias de media histórica, persistencia, AR(49) y vecino más próximo para que la comparación sea homogénea.

La fase genera tablas de selección de \(k\) por configuración, métricas de test, comparación compacta entre configuraciones y un resumen JSON. Las figuras muestran el RMSE de validación por \(k\), el RMSE de test por configuración, la comparación de métricas y una ventana real frente a predicha para las mejores configuraciones. La lectura de estos resultados se reserva para el capítulo de resultados; en este capítulo la fase queda definida como prueba de sensibilidad del procedimiento local.

\section{Comprobación sintética con mapa logístico}\label{sec:validacion-logistica}

La Fase 13 replica una parte sustancial del procedimiento sobre un sistema sintético conocido: el mapa logístico en régimen caótico. Esta fase no utiliza datos de Bitcoin. Su función es actuar como control metodológico: comprobar que las herramientas de reconstrucción, cuantificación y predicción local responden de forma razonable cuando la dinámica subyacente sí procede de un sistema determinista no lineal.

La serie limpia se genera mediante
\begin{equation}
    x_{t+1} = 4x_t(1-x_t),
\end{equation}
con \(x_0=0{,}123456789\). Se simulan 12.000 iteraciones y se descartan las primeras 1.000 como transitorio, quedando 11.000 observaciones útiles. A partir de la serie limpia se construyen dos variantes con ruido observacional gaussiano: una con ruido pequeño y otra con ruido moderado. La desviación típica del ruido se fija como proporción de la desviación típica de la serie limpia, y los valores resultantes se recortan al intervalo \([0,1]\).

El script \texttt{phase13\_\allowbreak logistic\_\allowbreak map\_\allowbreak validation.py} aplica a cada serie sintética la selección de retardo mediante AMI, la selección de dimensión mediante FNN y Cao, la reconstrucción por retardos, visualizaciones 2D y 3D, una estimación tipo Rosenstein y una evaluación predictiva local. En este caso el horizonte predictivo es 1, porque el mapa logístico define directamente \(x_{t+1}\) a partir de \(x_t\).

La partición temporal de cada serie sintética se realiza con un 60\% para entrenamiento, 20\% para validación y 20\% para test. La rejilla de vecinos es
\[
    k \in \{1,2,3,5,10,20,50\}.
\]
La selección de \(k\) se realiza en validación y la evaluación final se calcula en test. Se mantienen controles temporales análogos a los usados en Bitcoin: los vecinos de validación se buscan en entrenamiento y los de test en entrenamiento más validación.

La fase incluye también una referencia lineal autorregresiva seleccionada por BIC. A diferencia de Bitcoin, aquí se contempla explícitamente AR(0) como media histórica, además de AR(\(p\)) para \(p=1,\ldots,100\). Esta referencia sirve para comprobar si una aproximación lineal simple puede capturar la dinámica sintética o si el predictor local aprovecha mejor la estructura no lineal del sistema.

Las salidas incluyen CSV con las series sintéticas, tablas de AMI, FNN, Cao, parámetros seleccionados, curva de Rosenstein, métricas predictivas, selección AR por BIC y resumen JSON. Las figuras generadas muestran ejemplos de las series, curvas de AMI/FNN/Cao, embeddings 2D/3D, divergencia de Rosenstein, selección de \(k\), selección AR, mapa de transición de la serie limpia, predicciones real frente a predicho y comparación de métricas. La interpretación final de esta comprobación se desarrolla en el capítulo de resultados.

\section{Modelo HAR-logRV y exportación al MVP}\label{sec:har-logrv-exportacion}

La Fase 14 introduce un modelo HAR-logRV compacto como cierre predictivo operativo. El objetivo no es abrir una búsqueda amplia de modelos de volatilidad, sino incorporar una referencia específica para volatilidad realizada que sea simple, interpretable y exportable al MVP web.

El modelo utiliza tres escalas pasadas de volatilidad realizada:
\[
    \texttt{log\_\allowbreak rv\_\allowbreak past\_\allowbreak 12},\quad
    \texttt{log\_\allowbreak rv\_\allowbreak past\_\allowbreak 48},\quad
    \texttt{log\_\allowbreak rv\_\allowbreak past\_\allowbreak 288}.
\]
La variable objetivo es \texttt{log\_\allowbreak rv\_\allowbreak future\_\allowbreak 12}. La especificación implementada es
\begin{equation}
    \texttt{log\_\allowbreak rv\_\allowbreak future\_\allowbreak 12}
    =
    \beta_0
    + \beta_1 \texttt{log\_\allowbreak rv\_\allowbreak past\_\allowbreak 12}
    + \beta_2 \texttt{log\_\allowbreak rv\_\allowbreak past\_\allowbreak 48}
    + \beta_3 \texttt{log\_\allowbreak rv\_\allowbreak past\_\allowbreak 288}
    + \varepsilon_t.
\end{equation}

La estimación se realiza mediante mínimos cuadrados ordinarios sobre el tramo de entrenamiento. La implementación resuelve las ecuaciones normales y solo recurre a una regularización ridge muy pequeña si la matriz resulta numéricamente singular. Las tres variables del modelo se fijan antes de la evaluación, por lo que validación no se usa para selección de variables y test no interviene en entrenamiento ni en ajuste.

El script \texttt{phase14\_\allowbreak har\_\allowbreak logrv\_\allowbreak mvp.py} compara HAR-logRV con media histórica, persistencia, AR(49) a horizonte 12 y kNN local con \(\tau=137\), \(m=5\) y \(k=200\). Para que la comparación con kNN sea homogénea, se utiliza una muestra comparable de 5.000 puntos de test. Además, se calculan métricas sobre validación completa y test completo para los modelos que no requieren búsqueda de vecinos.

La fase mantiene una comprobación explícita de alineación temporal entre \texttt{log\_\allowbreak rv\_\allowbreak future\_\allowbreak 12(t)} y \texttt{log\_\allowbreak rv\_\allowbreak past\_\allowbreak 12(t+12)}. También documenta los controles anti-\textit{leakage}: HAR se entrena solo con train, las variables son pasadas, test no se usa para selección ni ajuste, \(k=200\) procede de la fase anterior y los vecinos de test del kNN se buscan solo en entrenamiento más validación.

Como salida principal, la fase exporta el artefacto \texttt{har\_\allowbreak logrv\_\allowbreak model.json} en \texttt{data/\allowbreak model\_\allowbreak artifacts/}. El artefacto contiene nombre y versión del modelo, fechas de entrenamiento, objetivo, variables, intercepto, coeficientes, tamaño de entrenamiento, posible regularización usada y métricas principales. Esta exportación permite que el MVP ejecute el modelo sin importar código del repositorio técnico.

Además del artefacto exportado, la fase genera tablas con coeficientes HAR, métricas de validación y test, predicciones de muestra, comparación de modelos, simulación del modo MVP con las últimas 1.000 velas y resumen JSON. Las figuras asociadas muestran comparación de métricas, real frente a predicho, distribución de errores absolutos, errores temporales, coeficientes y contexto de predicción para el modo MVP.

\section{Resultado final de la validación del modelo}\label{sec:resultado-validacion}

El procedimiento completo deja tres resultados técnicos diferenciados. En primer lugar, la cadena de análisis permite construir una serie de volatilidad realizada limpia, validada y sin fuga de información. En segundo lugar, las fases de caracterización, reconstrucción y cuantificación muestran estructura temporal no trivial en la volatilidad de Bitcoin, aunque sin evidencia suficiente para afirmar caos determinista. En tercer lugar, el bloque predictivo muestra que la volatilidad futura a una hora es parcialmente predecible frente a referencias simples.

El kNN local confirma que el espacio reconstruido contiene información útil: mejora a persistencia y aporta una lectura compatible con estados de volatilidad similares que tienden a tener comportamientos futuros parcialmente parecidos. Sin embargo, no supera de forma robusta al AR(49). Esto limita la interpretación no lineal y obliga a mantener el AR como benchmark fuerte.

La comprobación sintética con mapa logístico cumple una función distinta. En ese sistema, el kNN local supera claramente a la referencia lineal y la reconstrucción muestra una estructura más nítida. Por tanto, el pipeline responde ante una dinámica no lineal conocida. La diferencia entre ese resultado y Bitcoin refuerza la lectura prudente: en datos financieros reales, la señal aparece mezclada con ruido, heterocedasticidad y cambios de régimen.

El cierre predictivo lo proporciona HAR-logRV. Su mejora frente a AR(49) es marginal, pero combina buen rendimiento, simplicidad, bajo coste y facilidad de exportación. Por ello se adopta como modelo práctico principal del MVP. La reconstrucción no lineal y el kNN quedan como parte del análisis metodológico; HAR-logRV queda como modelo operativo recomendado.

En conjunto, la implementación cumple el objetivo técnico del capítulo: construir un procedimiento reproducible, validarlo en un sistema sintético, aplicarlo con cautela a Bitcoin y exportar un modelo práctico a una herramienta web autónoma.
